\chapter{الطوبولوجيا العامة}\label{ch:b3:topology}

أنجز مجلد السنة الجامعية 2 التحليلَ في الفضاءات المترية:
بالمسافات والكرات والمتتاليات. لكن المفاهيم الأساسية ---
الاتصال والتراص والترابط --- لا تذكر أبدًا القيمة العددية
لمسافة، بل عائلة \emph{\omterm{def:b3:topology:topology}{المجموعات المفتوحة}} التي تولّدها فحسب.
ويأخذ هذا الفصل تلك العائلة موضوعًا أوّليًا. والمكسب ليس التعميم
لذاته: فإنشاءات القسمة (الدائرة بوصفها $\R/\Z$، والفضاءات
الإسقاطية)، والجداءات، والطوبولوجيات الضعيفة في التحليل الدالي،
ليست أصلًا موضوعات مترية أولًا. نعيد بناء الاتصال، ثم نعالج
التراص بالتغطيات المفتوحة (مبرهنين على تكافئه مع تعريف السنة
الجامعية 2 بالمتتاليات في الفضاءات المترية)، ثم الترابط ---
ونختم بالمبرهنة القائلة إن تقابلًا \omterm{def:b3:topology:continuity}{متصلًا} لا يمكن أن يطابق $\R$
مع $\R^2$: \omterm{def:b3:topology:topology}{فالطوبولوجيا} قادرة على التمييز بين الأبعاد.

\section{الطوبولوجيات والمجموعات المفتوحة والاتصال}

\begin{definition}\label{def:b3:topology:topology}
\emph{الطوبولوجيا}\index{طوبولوجيا} على مجموعة $X$ هي عائلة
$\mathcal T$ من أجزاء $X$ --- تُسمّى \emph{المجموعات
المفتوحة}\index{مجموعة مفتوحة} --- بحيث: $\varnothing, X \in
\mathcal T$؛ وأي اتحاد
لمجموعات مفتوحة مفتوح؛ وأي تقاطع \emph{منتهٍ} لمجموعات مفتوحة
مفتوح. ويُسمّى الزوج $(X, \mathcal T)$ \emph{فضاءً طوبولوجيًا}. ومتممات
المجموعات المفتوحة \emph{مغلقة}. و\emph{الجوار}\index{جوار}
للنقطة $x$ مجموعةٌ تحتوي على مجموعة مفتوحة تحتوي على $x$.
\end{definition}

\begin{example}\label{ex:b3:topology:examples}
(a) الفضاء المتري، مع «المفتوح» كما في السنة الجامعية 2 (أي
اتحادات الكرات المفتوحة): وهي \emph{\omterm{def:b3:topology:topology}{الطوبولوجيا} المترية}؛ ويمكن
لمسافات مختلفة أن تعطي \omterm{def:b3:topology:topology}{الطوبولوجيا} نفسها (المسافات المتكافئة).
ويكون الفضاء \emph{قابلًا للتمتير} إذا نشأت طوبولوجيته عن مسافة
ما.
(b) \omterm{def:b3:topology:topology}{الطوبولوجيا} \emph{المتقطّعة} (وفيها جميع الأجزاء مفتوحة)
\omterm{def:b3:topology:topology}{والطوبولوجيا} \emph{الخشنة} $\{\varnothing, X\}$.
(c) \omterm{def:b3:topology:topology}{طوبولوجيا} \emph{المتممات المنتهية} على مجموعة لامنتهية:
المفتوح $=$ الخالي أو ذو المتمم المنتهي. وهي غير قابلة للتمتير،
كما سنرى (\cref{exo:b3:topology:3}).
(d) على $\R$، \omterm{def:b3:topology:topology}{الطوبولوجيا} المعتادة؛ وعلى $\bar\R = \R \cup
\{\pm\infty\}$، \omterm{def:b3:topology:topology}{طوبولوجيا}
الترتيب المولَّدة بالأشعة --- وهي تجعل «$x_n \to +\infty$» حالةً من
التقارب العادي.
\end{example}

\begin{definition}\label{def:b3:topology:interior}
من أجل $A \subseteq X$: يكون \emph{الداخل} $\mathring A$ أكبر
\omterm{def:b3:topology:topology}{مجموعة مفتوحة} داخل $A$ (أي اتحاد جميعها)؛ ويكون
\emph{الغلق}\index{غلق} $\bar A$ أصغر مجموعة مغلقة تحتوي على
$A$؛ وتكون \emph{الحافة} $\partial A = \bar A
\setminus \mathring A$. وتكون $A$
\emph{كثيفة}\index{مجموعة كثيفة} إذا كان $\bar A = X$. ويتحقق
$x \in \bar A$ إذا وفقط إذا لاقى كلُّ \omterm{def:b3:topology:topology}{جوار} للنقطة $x$ المجموعةَ
$A$ (إذ لو أخطأ \omterm{def:b3:topology:topology}{جوار} ما المجموعةَ $A$، لكان متمم لبّه المفتوح
مجموعة مغلقة أصغر تحيط بالمجموعة $A$؛ والعكس بالمثل).
\end{definition}

\begin{definition}\label{def:b3:topology:basis}
\emph{الأساس}\index{أساس طوبولوجيا} \omterm{def:b3:topology:topology}{للطوبولوجيا} $\mathcal T$ هو
عائلة $\mathcal B \subseteq \mathcal T$ بحيث تكون كل \omterm{def:b3:topology:topology}{مجموعة مفتوحة} اتحادًا لعناصر
$\mathcal B$ (مثل الكرات المفتوحة في فضاء متري؛ والفترات
المفتوحة في $\R$). وتكون عائلة $\mathcal B$ من أجزاء $X$ أساسًا
\omterm{def:b3:topology:topology}{لطوبولوجيا} \emph{ما} إذا وفقط إذا غطّت $X$ وكان، من أجل $B_1, B_2 \in \mathcal B$
و $x \in B_1 \cap B_2$، عنصرٌ $B_3 \in \mathcal B$ يحقق $x \in B_3 \subseteq B_1\cap B_2$ --- وعندئذٍ تكون «اتحادات
العناصر» \omterm{def:b3:topology:topology}{طوبولوجيا}، وهي \omterm{def:b3:topology:topology}{الطوبولوجيا} \emph{المولَّدة}
بالعائلة $\mathcal B$.
\end{definition}

\begin{definition}\label{def:b3:topology:continuity}
يكون $f \colon X \to Y$ \emph{متصلًا}\index{تطبيق متصل} إذا كان $f^{-1}(V)$
مفتوحًا من أجل كل مفتوح $V \subseteq Y$ --- وهذا يكافئ أن تكون
سوابق المجموعات المغلقة مغلقة؛ ويكافئ أن يكون $f^{-1}(V)$
جوارًا للنقطة $x$ من أجل كل $x$ وكل \omterm{def:b3:topology:topology}{جوار} $V$ للنقطة $f(x)$ (أي
الاتصال \emph{عند} كل $x$). ويكفي التحقق على أساس للفضاء $Y$.
وتركيب التطبيقات المتصلة متصل. و\emph{التماثل
الطوبولوجي}\index{تماثل طوبولوجي} هو تقابل متصل مقلوبه متصل؛
وتدرس \omterm{def:b3:topology:topology}{الطوبولوجيا} الخصائصَ التي تحفظها التماثلات الطوبولوجية.
\end{definition}

\begin{proposition}[الاتصال والغلق؛ اللصق]
\label{prop:b3:topology:gluing}
(a) يكون $f$ \omterm{def:b3:topology:continuity}{متصلًا} إذا وفقط إذا كان $f(\bar A) \subseteq \overline{f(A)}$ من أجل كل
$A \subseteq X$.
(b) إذا كان $X = F_1 \cup F_2$ حيث $F_1, F_2$ مغلقتان وكان قصر $f\colon X
\to Y$ على
كل $F_i$ \omterm{def:b3:topology:continuity}{متصلًا}، فإن $f$ متصل.
\end{proposition}

\begin{proof}
(a) إذا كان $f$ \omterm{def:b3:topology:continuity}{متصلًا}: فإن $f^{-1}(\overline{f(A)})$ مغلقة وتحتوي على $A$، ومنه
تحتوي على $\bar A$. وبالعكس، نطبّق المحك على $A = f^{-1}(F)$
حيث $F$ مغلقة: $f(\bar A) \subseteq
\overline{f(f^{-1}(F))} \subseteq \bar F = F$، ومنه $\bar A
\subseteq f^{-1}(F) = A$: أي إن سوابق المجموعات
المغلقة مغلقة.
(b) من أجل $F \subseteq Y$ مغلقة: $f^{-1}(F) =
(f\restriction_{F_1})^{-1}(F) \cup (f\restriction_{F_2})^{-1}
(F)$، وهو اتحاد مجموعتين
مغلقتين في $F_1$ و $F_2$ على الترتيب، ومنه فهو مغلق في $X$ (لأن
$F_i$ مغلقة: فالمغلق داخل مغلق مغلقٌ).
\end{proof}

\begin{definition}\label{def:b3:topology:hausdorff}
يكون $X$ \emph{هاوسدورفيًا}\index{فضاء هاوسدورف} (أو
\emph{مفصولًا}) إذا كان لكل نقطتين متمايزتين جواران منفصلان.
والفضاءات المترية هاوسدورفية (بكرات نصف قطرها $d(x,y)/2$).
وتُقال متتالية $(x_n)$ إنها \emph{تتقارب} إلى $x$ إذا احتوى كل
\omterm{def:b3:topology:topology}{جوار} للنقطة $x$ على جميع الحدود $x_n$ عدا عددًا منتهيًا منها؛
وفي فضاء هاوسدورفي تكون النهايات وحيدة (إذ سيكون لنهايتين
جواران منفصلان يحتوي كلٌّ منهما على ذيل). وفي فضاء هاوسدورفي
تكون النقط --- ومنه المجموعات المنتهية --- مغلقة.
\end{definition}

\begin{remark}\label{rem:b3:topology:sequences}
في الفضاءات المترية، تكشف المتتاليات كل شيء: فيتحقق
$x \in \bar A$ إذا وفقط إذا تقاربت متتالية من $A$ إلى $x$ (بأخذ
$x_n \in A \cap
B(x, 1/n)$)، ويكون $f$ \omterm{def:b3:topology:continuity}{متصلًا} إذا وفقط إذا كان \omterm{def:b3:topology:continuity}{متصلًا} بالمتتاليات.
أما في الفضاءات العامة فيخفق التكافؤان كلاهما؛ وتنتمي بدائل
المتتاليات الصحيحة (المرشّحات والشبكات) إلى مقرّر أكثر تقدّمًا.
وسنذكر النتائج القائمة على المتتاليات في الفضاءات المترية
والنتائج القائمة على التغطيات في العموم --- ونبرهن على تكافئها
حيث تتكافأ.
\end{remark}

\section{الفضاءات الجزئية والجداءات وحواصل القسمة}

\begin{definition}\label{def:b3:topology:constructions}
ثلاث طرائق لصنع فضاءات جديدة من قديمة:\index{طوبولوجيا
جدائية}\index{طوبولوجيا القسمة}
\begin{enumerate}
  \item \emph{الفضاء الجزئي}: على $A \subseteq X$، تكون
        \omterm{def:b3:topology:topology}{المجموعات المفتوحة} هي $U \cap A$ حيث $U$ مفتوحة في $X$
        --- وهي أخشن \omterm{def:b3:topology:topology}{طوبولوجيا} تجعل الاحتواء \omterm{def:b3:topology:continuity}{متصلًا}.
  \item \emph{الجداء}: على $X \times Y$ (وعلى الجداءات
        المنتهية)، \omterm{def:b3:topology:topology}{الطوبولوجيا} التي أساسها \emph{الصناديق
        المفتوحة} $U \times
        V$؛ وعلى جداء لامنتهٍ $\prod_i X_i$،
        يتألف \omterm{def:b3:topology:basis}{الأساس} من الصناديق $\prod U_i$ التي يكون فيها
        $U_i = X_i$ من أجل \emph{جميع الأدلّة عدا عددًا منتهيًا}
        منها --- وهي أخشن \omterm{def:b3:topology:topology}{طوبولوجيا} تجعل كل إسقاط \omterm{def:b3:topology:continuity}{متصلًا}.
  \item \emph{القسمة}: إذا كان $\sim$ علاقة تكافؤ على $X$ وكان
        $\pi\colon X \to X/{\sim}$ الإسقاطَ، نقول إن $V
        \subseteq X/{\sim}$ مفتوحة إذا وفقط إذا كانت
        $\pi^{-1}(V)$ مفتوحة --- وهي أدقّ \omterm{def:b3:topology:topology}{طوبولوجيا} تجعل $\pi$
        \omterm{def:b3:topology:continuity}{متصلًا}.
\end{enumerate}
\end{definition}

\begin{proposition}[الخصائص الشمولية]
\label{prop:b3:topology:universal}
(a) يكون $f \colon Z \to X \times Y$ \omterm{def:b3:topology:continuity}{متصلًا} إذا وفقط إذا كانت المركّبتان
$\pi_X \circ f$ و $\pi_Y \circ f$ متصلتين (والأمر نفسه من أجل
جداءات كيفية).
(b) يكون $g \colon X/{\sim} \to Z$ \omterm{def:b3:topology:continuity}{متصلًا} إذا وفقط إذا كان $g \circ \pi
\colon X \to Z$ \omterm{def:b3:topology:continuity}{متصلًا}.
\end{proposition}

\begin{proof}
(a) اللزوم: بالتركيب. الكفاية: يكفي التحقق من سوابق صناديق
\omterm{def:b3:topology:basis}{الأساس}: $f^{-1}(U \times V) = (\pi_Xf)^{-1}(U)
\cap (\pi_Yf)^{-1}(V)$، وهي مفتوحة (لأن عدد العوامل $\neq X_i$ منتهٍ في
الحالة اللامنتهية). (b) اللزوم: بالتركيب. الكفاية: من أجل
$W \subseteq Z$ مفتوحة، تكون $\pi^{-1}(g^{-1}(W)) = (g\pi)^{-1}(W)$ مفتوحة، وهذا يعني بحكم تعريف
\omterm{def:b3:topology:constructions}{طوبولوجيا القسمة} أن $g^{-1}(W)$ مفتوحة.
\end{proof}

\begin{example}\label{ex:b3:topology:circle}
حاصل القسمة $\R/\Z$ (بمطابقة $x$ مع $x + n$) مماثل طوبولوجيًا
للدائرة $S^1 = \{z \in \C : \abs z = 1\}$: إذ ينتقل التطبيق $x
\mapsto \eu^{2\iu\pi x}$ إلى تقابل متصل
$\R/\Z \to S^1$ (\cref{prop:b3:topology:universal}(b))؛ ومقلوبه متصل بحجّة التراص في
\cref{cor:b3:topology:compacthomeo} أدناه (ويفصّل \cref{exo:b3:topology:5} كل شيء، بما في ذلك سبب كون
$\R/\Z$ \omterm{def:b3:topology:hausdorff}{هاوسدورفيًا} \omterm{def:b3:topology:compact}{ومتراصًّا}). وبالمثل فإن $[0,1]$ بلصق طرفيه
هو $S^1$، والمربع بلصق أضلاعه المتقابلة هو الطارة، ويصير اللصق
أخيرًا مبرهنةً لا صورة.
\end{example}

\section{التراص}

\begin{definition}\label{def:b3:topology:compact}
\emph{التغطية المفتوحة} للفضاء $X$ هي عائلة $(U_i)_{i\in I}$ من
\omterm{def:b3:topology:topology}{المجموعات المفتوحة} تحقق $\bigcup U_i = X$. ويكون $X$
\emph{متراصًّا}\index{فضاء متراص} إذا كان \omterm{def:b3:topology:hausdorff}{هاوسدورفيًا} وقبلت كل
تغطية مفتوحة \emph{تغطية جزئية منتهية}. وهذا يكافئ (بأخذ
المتممات): أن يكون لكل عائلة من المجموعات المغلقة تحقق
\emph{خاصية التقاطع المنتهي} (أي إن لكل عائلة جزئية منتهية
تقاطعًا غير خالٍ) تقاطعٌ كلي غير خالٍ.
\end{definition}

\begin{theorem}[الخصائص الأولى]\label{thm:b3:topology:compactprops}
ليكن $X$ \omterm{def:b3:topology:compact}{متراصًّا}.
\begin{enumerate}
  \item كل جزء مغلق من $X$ متراص؛ وكل جزء متراص من فضاء
        هاوسدورفي مغلق.
  \item صورة \omterm{def:b3:topology:compact}{فضاء متراص} \omterm{def:b3:topology:continuity}{بتطبيق متصل} داخل فضاء هاوسدورفي متراصةٌ.
        وعلى وجه الخصوص، يكون كل \omterm{def:b3:topology:continuity}{تطبيق متصل} $f\colon X
        \to \R$ محدودًا
        ويبلغ حدّيه.
  \item لكل متتالية متناقصة من الأجزاء المغلقة غير الخالية من
        $X$ تقاطعٌ غير خالٍ.
\end{enumerate}
\end{theorem}

\begin{proof}
(1) لتكن $F \subseteq X$ مغلقة ولتكن $(U_i)$ تغطية مفتوحة
للمجموعة $F$ (بمفتوحات من $X$): فبإضافة $X \setminus F$ نحصل على تغطية
مفتوحة للفضاء $X$؛ وتغطية جزئية منتهية منها، منقوصًا منها
$X \setminus F$، تغطّي $F$. وكون الفضاء الجزئي \omterm{def:b3:topology:hausdorff}{هاوسدورفيًا}
واضح. وبالعكس، لتكن $K \subseteq Y$ متراصة و $Y$ \omterm{def:b3:topology:hausdorff}{هاوسدورفيًا}
وليكن $y \notin K$: فمن أجل كل $x \in K$ نختار مفتوحتين
منفصلتين $U_x \ni x$ و $V_x \ni y$؛ ويغطّي عدد منتهٍ من
المجموعات $U_x$ المجموعةَ $K$، ويكون تقاطع المجموعات $V_x$
الموافقة جوارًا للنقطة $y$ منفصلًا عن تغطية $K$، ومنه عن $K$:
أي إن متمم $K$ مفتوح.

(2) إذا غطّت $(V_j)$ المجموعةَ $f(X)$، غطّت $(f^{-1}(V_j))$
الفضاءَ $X$؛ وتأتي تغطية جزئية منتهية في الأسفل من التغطية
الجزئية المنتهية في الأعلى. والصورة هاوسدورفية بوصفها فضاءً
جزئيًا. ومن أجل $f$ حقيقي: تكون $f(X)$ متراصة في $\R$، ومنه
مغلقة ومحدودة (بالتغطية بالفترات $(-n, n)$ من أجل الحد؛
ومغلقة حسب (1))، وتحتوي المجموعة المغلقة المحدودة على حدّها
الأعلى.

(3) إذا كان $\bigcap F_n = \varnothing$، فإن المفتوحات $X \setminus F_n$ تغطّي $X$؛ ويكفي عدد
منتهٍ منها، ومنه $F_{n_0} =
\varnothing$ (لأن المتتالية متناقصة) --- وهو تناقض.
\end{proof}

\begin{corollary}\label{cor:b3:topology:compacthomeo}
كل \emph{تقابل} متصل من \omterm{def:b3:topology:compact}{فضاء متراص} إلى فضاء هاوسدورفي هو \omterm{def:b3:topology:continuity}{تماثل
طوبولوجي}.
\end{corollary}

\begin{proof}
يكون المقلوب \omterm{def:b3:topology:continuity}{متصلًا} إذا وفقط إذا كانت الصور المباشرة للمجموعات
المغلقة مغلقة؛ وكل مغلقة $F$ متراصة (\cref{thm:b3:topology:compactprops}(1))، وصورتها
متراصة (2)، ومنه فهي مغلقة (1) في المستقر الهاوسدورفي.
\end{proof}

\begin{theorem}[الجداءات المنتهية]\label{thm:b3:topology:tychonoff}
كل جداء منتهٍ لفضاءات متراصة متراصٌّ.\index{مبرهنة تيخونوف
(الحالة المنتهية)}
\end{theorem}

\begin{proof}
يكفي أن نعالج $X \times Y$. وصفة هاوسدورف موروثة (بالفصل في
إحدى الإحداثيتين). ولتكن $(W_i)$ تغطية مفتوحة للفضاء $X
\times Y$؛
ويمكننا أن نفترض أن المجموعات $W_i$ صناديق أساسية $U_i \times
V_i$
(بالتنقيح: فكل نقطة تقع في صندوق داخل $W_i$ ما؛ وتعطي تغطية
جزئية منتهية بالصناديق تغطيةً جزئية منتهية بالمجموعات $W_i$).
لنثبّت $x \in X$: تكون \emph{الشريحة} $\{x\}\times Y \cong Y$ متراصة، ومنه
يغطّيها عدد منتهٍ من الصناديق $U_1\times V_1, \dots, U_k\times V_k$، مع $x \in U_j$ من أجل كل
$j$؛ وعندئذٍ تكون $U^x = \bigcap_{j \leq k} U_j$ جوارًا مفتوحًا للنقطة $x$ بحيث يُغطّى
$U^x \times Y$ بعدد منتهٍ من الصناديق (وهي \emph{مبرهنة الأنبوب
المساعدة}\index{مبرهنة الأنبوب المساعدة}: فمن أجل $(x', y) \in U^x \times Y$، يكون
$y \in V_j$ من أجل $j$ ما لأن الصناديق غطّت $\{x\}\times Y$ عند
المستوى $y$، ولدينا $x' \in U^x \subseteq U_j$، ومنه $(x', y) \in U_j
\times V_j$). والآن يغطّي عدد منتهٍ
من المجموعات $U^x$ الفضاءَ المتراص $X$؛ وتغطّي مجموعات الصناديق
المنتهية الموافقة الفضاءَ $X \times Y$.
\end{proof}

\begin{theorem}[التراص في الفضاءات المترية]
\label{thm:b3:topology:metriccompact}
من أجل فضاء متري $(X, d)$، تتكافأ العبارات التالية:\index{التراص
بالمتتاليات}
\begin{enumerate}
  \item الفضاء $X$ متراص (بوريل--لوبيغ)؛
  \item لكل متتالية في $X$ متتالية جزئية متقاربة (وهو \omterm{thm:b3:topology:metriccompact}{التراص
        بالمتتاليات} --- أي تعريف السنة الجامعية 2)؛
  \item الفضاء $X$ تام و\emph{محدود كليًا}: أي إن عددًا منتهيًا
        من الكرات ذات نصف القطر $\varepsilon$ يغطّي $X$ من أجل
        كل $\varepsilon > 0$.
\end{enumerate}
\end{theorem}

\begin{proof}
(1)$\Rightarrow$(2): لتكن $(x_n)$ بلا متتالية جزئية متقاربة.
عندئذٍ يكون لكل $x \in X$ كرة مفتوحة $B_x$ لا تحتوي على $x_n$
إلا من أجل عدد منتهٍ من الأدلّة $n$ (وإلا تقاربت متتالية جزئية
إلى $x$: بأخذ أنصاف الأقطار $1/k$). ويغطّي عدد منتهٍ من الكرات
$B_x$ الفضاءَ $X$، ومنه لا يوجد سوى عدد منتهٍ من الأدلّة $n$:
وهو سخف.

(2)$\Rightarrow$(3): التمام: فمتتالية كوشي لها متتالية جزئية
متقاربة تتقارب (السنة الجامعية 2). والحد الكلي: إذا لم يقبل
$\varepsilon$ ما تغطية منتهية، اخترنا بالتتابع $x_{n+1}$ خارج
$\bigcup_{k\leq n} B(x_k, \varepsilon)$: فتحقق المتتالية $d(x_m, x_n) \geq \varepsilon$ من أجل $m \neq n$، ولا تكون لها
متتالية جزئية كوشي، ولا متقاربة.

(3)$\Rightarrow$(1): أولًا، تستلزم (3) العبارةَ (2): إذ إذا
أُعطيت $(x_n)$، غطّينا $X$ بعدد منتهٍ من الكرات ذات نصف القطر
$1$: فتحتوي إحداها، ولتكن $B_1$، على متتالية جزئية؛ ثم نغطّي
$X$ بكرات نصف قطرها $1/2$: فتحتوي إحداها على متتالية جزئية
أخرى؛ ونكرّر ونقطّر: فتكون المتتالية الجزئية القطرية متتالية
كوشي (لأن حدّين بعد المرحلة $k$ يقعان في كرة مشتركة نصف قطرها
$2^{-k}$، إلى حدود العامل $2\cdot$ المعتاد)، ومنه فهي متقاربة.
والآن لتكن $(U_i)$ تغطية مفتوحة ولنفترض عدم وجود تغطية جزئية
منتهية. \emph{حجّة عدد لوبيغ}: من أجل كل $n$، توجد كرة
$B(y_n, 2^{-n})$ لا يغطّيها عدد منتهٍ من المجموعات $U_i$ ---
وبالفعل، نغطّي $X$ بعدد منتهٍ من الكرات ذات نصف القطر $2^{-n}$؛
فلو غُطّيت كلٌّ منها بعدد منتهٍ لغُطّي $X$ كذلك. وحسب (2)، توجد
متتالية جزئية $y_{n_k} \to y$؛ نختار $i$ يحقق $y \in U_i$
و $r > 0$ يحقق $B(y, r) \subseteq U_i$. ومن أجل $k$ كبير، $B(y_{n_k}, 2^{-n_k}) \subseteq B(y, r) \subseteq U_i$: أي مغطّاة
بمجموعة \emph{واحدة} $U_i$ --- وهو تناقض.
\end{proof}

\begin{corollary}[هاين--بوريل؛ هاين]
\label{cor:b3:topology:heineborel}
(a) يكون جزء من $\R^n$ \omterm{def:b3:topology:compact}{متراصًّا} إذا وفقط إذا كان مغلقًا
ومحدودًا.
(b) كل \omterm{def:b3:topology:continuity}{تطبيق متصل} من فضاء متري متراص إلى فضاء متري متصلٌ
بانتظام.\index{مبرهنة هاين--بوريل}\index{مبرهنة هاين}
\end{corollary}

\begin{proof}
(a) مغلق ومحدود $\Rightarrow$ محتوًى في مكعب $[-M,M]^n$، وهو
متراص: إذ إن $[-M, M]$ متراص (بالمتتاليات، ببولتزانو--فايرشتراس
--- أو مباشرةً بالتنصيف من أجل التغطيات)، وتعالج \cref{thm:b3:topology:tychonoff}
الجداءَ؛ ثم نطبّق \cref{thm:b3:topology:compactprops}(1). وبالعكس، يكون الجزء المتراص
مغلقًا (\cref{thm:b3:topology:compactprops}(1)) ومحدودًا (بالتغطية بكرات متحدة المركز).

(b) ليكن $f \colon X \to Y$ و $\varepsilon > 0$. تغطّي الكرات $B\bigl(x, \delta_x\bigr)$ مع $f\bigl(B(x,
2\delta_x)\bigr)\subseteq B\bigl(f(x), \varepsilon/2\bigr)$ الفضاءَ
$X$؛ نستخرج تغطية جزئية منتهية $B(x_j, \delta_{x_j})$ ونضع $\delta = \min_j \delta_{x_j}$. فإذا كان
$d(x, x') < \delta$: فإن $x
\in B(x_j, \delta_{x_j})$ من أجل $j$ ما، وعندئذٍ يكون كلٌّ من $x, x'
\in B(x_j, 2\delta_{x_j})$،
ومنه $d(f(x), f(x')) < \varepsilon$.
\end{proof}

\begin{definition}\label{def:b3:topology:locallycompact}
يكون $X$ \emph{متراصًّا محليًا}\index{فضاء متراص محليا} إذا كان
\omterm{def:b3:topology:hausdorff}{هاوسدورفيًا} وكان لكل نقطة \omterm{def:b3:topology:topology}{جوار} متراص (مثل $\R^n$؛ والأجزاء
المفتوحة من $\R^n$؛ والفضاءات المتقطّعة --- لكن ليس $\Q$، انظر
\cref{exo:b3:topology:9}). وكل \omterm{def:b3:topology:compact}{فضاء متراص} محليًا ينغمر في \omterm{def:b3:topology:compact}{فضاء متراص}: وهو
\emph{الترصيص بنقطة واحدة}\index{الترصيص بنقطة واحدة} $\hat X = X
\cup \{\infty\}$،
ومفتوحاته هي مفتوحات $X$ مضافًا إليها متممات (في $\hat X$)
الأجزاء المتراصة من $X$. ونتحقق من البديهيات مباشرةً؛ ويكون
$\hat X$ \omterm{def:b3:topology:compact}{متراصًّا} (إذ لكل تغطية عنصرٌ يحتوي على $\infty$،
ومتممه متراص ويغطّيه عدد منتهٍ من العناصر الأخرى) \omterm{def:b3:topology:hausdorff}{وهاوسدورفيًا}
(بفصل $x$ عن $\infty$ بجوار متراص للنقطة $x$ ومتممِه). مثال:
$\hat\R \cong S^1$، و $\widehat{\R^n} \cong S^n$ بالإسقاط المجسّم (\cref{exo:b3:topology:11}).
\end{definition}

\section{الترابط}

\begin{definition}\label{def:b3:topology:connected}
يكون $X$ \emph{مترابطًا}\index{فضاء مترابط} إذا لم يكن اتحاد
مجموعتين مفتوحتين منفصلتين غير خاليتين --- وهذا يكافئ أن تكون
أجزاؤه المفتوحة والمغلقة معًا هي $\varnothing$ و $X$ لا غير؛
ويكافئ أن يكون كل \omterm{def:b3:topology:continuity}{تطبيق متصل} $X \to \{0, 1\}$ (متقطّع) ثابتًا.
ويكون جزء مترابطًا إذا كان كذلك بوصفه فضاءً جزئيًا.
\end{definition}

\begin{theorem}\label{thm:b3:topology:connectedbasics}
\begin{enumerate}
  \item أجزاء $\R$ المترابطة هي بالضبط الفترات.
  \item صور المجموعات المترابطة بتطبيقات متصلة مترابطةٌ (ومن هنا
        مبرهنة القيم الوسطى: فصورة تطبيق حقيقي متصل على \omterm{def:b3:topology:connected}{فضاء
        مترابط} فترةٌ).
  \item إذا كانت $(A_i)$ مترابطة ولها نقطة مشتركة، فإن $\bigcup
        A_i$
        مترابطة. وإذا كانت $A$ مترابطة وكان $A \subseteq
        B \subseteq \bar A$، فإن $B$
        مترابطة.
  \item الجداءات المنتهية لفضاءات مترابطة مترابطةٌ.
\end{enumerate}
\end{theorem}

\begin{proof}
نستعمل في كل ما يلي محك $\{0,1\}$: أي إن $f \colon X \to
\{0,1\}$ المتصل يجب أن
يكون ثابتًا.

(1) كل مجموعة $A$ ليست فترةً تُخطئ نقطة $c$ بين $a, b \in A$:
$A = (A \cap (-\infty, c)) \sqcup (A \cap (c, +\infty))$
فتفصلها. وبالعكس، لتكن $I$ فترةً وليكن $f \colon I
\to \{0,1\}$ \omterm{def:b3:topology:continuity}{متصلًا} مع
$f(a) = 0$ و $f(b) = 1$ و $a < b$. نضع $c = \sup\{x \in [a,b] : f(x) = 0\}$؛ فيفرض الاتصال
عند $c$ أن $f(c) = 0$ (لأنه نهاية قيم $0$: إذ إن كل \omterm{def:b3:topology:topology}{جوار}
للنقطة $c$ يلاقي $\{f = 0\}$، و $\{f=0\}$ مغلقة)، وعندئذٍ يكون
$c < b$ مع $f \equiv 1$ على $(c, b]$، ومنه $f(c) = 1$ بحجّة
\omterm{def:b3:topology:interior}{الغلق} نفسها على $\{f = 1\} \ni c$: وهو تناقض.

(2) يعطي \omterm{def:b3:topology:continuity}{التطبيق المتصل} $g \colon f(X) \to \{0,1\}$ أن $g \circ f$ ثابت، ومنه فإن
$g$ ثابت على الصورة.

(3) يكون \omterm{def:b3:topology:continuity}{التطبيق المتصل} $f\colon \bigcup A_i \to \{0,1\}$ ثابتًا على كل $A_i$، بالقيمة
نفسها عند النقطة المشتركة. وأما من أجل \omterm{def:b3:topology:interior}{الغلق}: فإن $f \colon B \to \{0,1\}$ ثابت
$= c$ على $A$؛ وكل $x \in B \subseteq \bar A$ يقع في \omterm{def:b3:topology:interior}{غلق} $A$، وتكون $f^{-1}
(f(x))$ جوارًا
للنقطة $x$ (بوصفها سابق مفتوحة)، ولا بد أن تلاقي $A$: ومنه
$f(x) = c$.

(4) من أجل $X \times Y$ ومن أجل $f\colon X\times Y \to \{0,1\}$: يُوصَل أي نقطتين
$(x,y), (x',y')$ بواسطة «الكوع» $\{x\}\times Y \cup X \times \{y'\}$، وهو اتحاد مجموعتين
مترابطتين (مماثلتين طوبولوجيًا للفضاءين $Y$ و $X$) تتلاقيان عند
$(x, y')$: ومنه، حسب (3) و (2)، يتوافق $f$ على النقطتين.
\end{proof}

\begin{definition}\label{def:b3:topology:pathconnected}
يكون $X$ \emph{مترابطًا بالمسارات}\index{فضاء مترابط بالمسارات}
إذا كان أي نقطتين موصولتين بواسطة \emph{مسار} (أي \omterm{def:b3:topology:continuity}{تطبيق متصل} $\gamma
\colon [0,1] \to X$).
والترابط بالمسارات يستلزم الترابط: إذ إن قيمتَي \omterm{def:b3:topology:continuity}{تطبيق متصل}
$f \colon X \to \{0,1\}$ عند $x, y$ هما قيمتان للتطبيق الثابت $f\circ\gamma$
(\cref{thm:b3:topology:connectedbasics}(1)--(2)). والأجزاء المحدّبة من الفضاءات المعيارية
مترابطة بالمسارات (بالقطع)؛ وكذلك $\R^n
\setminus \{0\}$ من أجل $n \geq 2$
(بالدوران حول المبدأ)، و $S^n$ من أجل $n \geq 1$ (بإسقاط
المسارات من $\R^{n+1}\setminus
\{0\}$).
\end{definition}

\begin{example}[منحني جيب الطوبولوجي]
\label{ex:b3:topology:sinecurve}
لتكن $\Gamma = \{(x, \sin\frac1x) : 0 < x \leq 1\}$ ولتكن $S =
\bar\Gamma = \Gamma \cup (\{0\}\times[-1,1])$ (فكل نقطة $(0,
y)$ مع $\abs y \leq 1$
نهايةٌ لنقط من $\Gamma$: بحل $\sin\frac1x = y$ قرب $0$). عندئذٍ تكون $S$
\emph{مترابطة} --- لأنها \omterm{def:b3:topology:interior}{غلق} المجموعة المترابطة $\Gamma$، وهي
صورة متصلة للفترة $(0, 1]$ (\cref{thm:b3:topology:connectedbasics}(3)) --- لكنها
\emph{غير مترابطة بالمسارات}: إذ إن مسارًا من $(1, \sin 1)$ إلى
$(0,0)$ سيضطر إلى قطع فواصل $\to 0$ بينما يتذبذب الترتيب بين
$\pm1$؛ ويجعل \cref{exo:b3:topology:9} ذلك دقيقًا. فالترابط والترابط
بالمسارات مختلفان فعلًا.\index{منحني جيب الطوبولوجي}
\end{example}

\begin{omfigure}
\begin{tikzpicture}
\begin{axis}[omaxis, width=10.5cm, height=4.6cm,
    xmin=-0.02, xmax=0.55, ymin=-1.25, ymax=1.25,
    xtick={0, 0.25, 0.5}, ytick={-1, 1}]
  \addplot[omThm, thick, domain=0.006:0.55, samples=1200]
    {sin(deg(1/x))};
  \addplot[omDef, very thick] coordinates {(0,-1) (0,1)};
\end{axis}
\end{tikzpicture}

{\small \omterm{ex:b3:topology:sinecurve}{منحني جيب الطوبولوجي}: يتراكم منحني $\sin\frac1x$
(بالأحمر) على القطعة $\{0\}\times[-1,1]$ بأكملها (بالأزرق). والاتحاد مترابط
لكنه غير مترابط بالمسارات: فلا يمكن لأي \omterm{def:b3:topology:pathconnected}{مسار} متصل أن يعبر
التذبذبات اللامنتهية العدد في زمن وسيطي منتهٍ.}
\end{omfigure}

\begin{definition}\label{def:b3:topology:components}
\emph{المركّبة المترابطة}\index{مركبة مترابطة} للنقطة $x
\in X$ هي
اتحاد جميع الأجزاء المترابطة الحاوية على $x$ --- وهي أكبرها
(\cref{thm:b3:topology:connectedbasics}(3)). وتجزّئ المركّبات الفضاءَ $X$ وتكون مغلقة (لأن
غلوق المجموعات المترابطة مترابطة). ويكون $X$ \emph{غير مترابط
كليًا} إذا كانت جميع مركّباته أحادية (مثل $\Q$؛ ومجموعة كانتور
في مسألة نهاية الأسبوع).
\end{definition}

\begin{theorem}\label{thm:b3:topology:rnotr2}
الفضاء $\R$ ليس مماثلًا طوبولوجيًا لأي $\R^n$ مع $n \geq 2$.
\end{theorem}

\begin{proof}
لنفترض أن $h \colon \R \to \R^n$ \omterm{def:b3:topology:continuity}{تماثل طوبولوجي}. وبنزع نقطة: يكون $\R \setminus \{0\}$
مماثلًا طوبولوجيًا للفضاء $\R^n \setminus
\{h(0)\}$. لكن $\R\setminus\{0\}$ غير مترابط، بينما
$\R^n\setminus\{\text{نقطة}\}$ مترابط بالمسارات من أجل $n \geq
2$ (\cref{def:b3:topology:pathconnected}؛ بنقل
النقطة إلى $0$): والترابط صامدٌ بالتماثلات الطوبولوجية
(\cref{thm:b3:topology:connectedbasics}(2)) --- وهو تناقض. (وأما أن $\R^2 \not\cong \R^3$ فيتطلب صوامد
أدقّ --- من \omterm{def:b3:topology:topology}{الطوبولوجيا} الجبرية؛ وتُظهر مسألة نهاية الأسبوع
الخطر: إذ إن \emph{التطبيقات الشاملة المتصلة} $\R \to \R^2$ موجودة
فعلًا.)
\end{proof}

\begin{method}\label{met:b3:topology:connectargument}
لإثبات أن مجموعة مترابطة: أظهرها بوصفها صورة متصلة، أو اتحاد
مجموعات مترابطة متراكبة، أو غلقًا، أو جداءً (\cref{thm:b3:topology:connectedbasics})؛ ومن
أجل أجزاء الفضاءات المعيارية، برهن على الترابط بالمسارات بمسارات
صريحة (قطع، أقواس، أكواع). ولإثبات أن فضاءين \emph{غير}
مماثلين طوبولوجيًا: جد صامدًا طوبولوجيًا يختلف --- كالتراص، أو
الترابط، أو عدد المركّبات، أو المركّبات بعد حذف مجموعة منتهية
مختارة بعناية (وهي حيلة \cref{thm:b3:topology:rnotr2}: وهي تبرهن أيضًا على $[0,1) \not\cong
(0,1)$
وعلى $S^1 \not\cong [0,1]$).
\end{method}

\section{تمارين}

\begin{exercise}[$\star$]\label{exo:b3:topology:1}
(a) اسرد جميع الطوبولوجيات على $\{a, b\}$، وصنّفها إلى حدود
\omterm{def:b3:topology:continuity}{التماثل الطوبولوجي}، وعيّن أيها مترابط وأيها هاوسدورفي.
(b) في $\R$ (\omterm{def:b3:topology:topology}{بالطوبولوجيا} المعتادة)، احسب داخل $\Q$ وغلقه
وحافته، وكذلك من أجل $[0,1) \cup \{2\}$ ومن أجل $\{1/n : n \geq
1\}$.
\end{exercise}

\begin{exercise}[$\star$]\label{exo:b3:topology:2}
(a) برهن على أن $f\colon X \to Y$ متصل إذا وفقط إذا كانت $f^{-1}(B)$
مفتوحة من أجل كل $B$ من أساس مثبَّت للفضاء $Y$.
(b) برهن على أن $f = \mathbf 1_{[0,\infty)} \colon \R \to \R$ ليس \omterm{def:b3:topology:continuity}{متصلًا} لكنه \emph{متصل من اليمين}.
وتحقق من أن الفترات نصف المفتوحة $[a, b)$ تشكّل أساسًا
\omterm{def:b3:topology:topology}{لطوبولوجيا} على المنطلق $\R$ (وهو \emph{مستقيم سورغنفراي})، ومن
أن هذه \omterm{def:b3:topology:topology}{الطوبولوجيا} أدقّ تمامًا من المعتادة، ومن أن الدالة
$\R \to
\R$ (بالمستقر المعتاد) تكون متصلة من مستقيم سورغنفراي إذا
وفقط إذا كانت متصلة من اليمين عند كل نقطة.
\end{exercise}

\begin{exercise}[$\star$]\label{exo:b3:topology:3}
على مجموعة لامنتهية مزوّدة \omterm{def:b3:topology:topology}{بطوبولوجيا} المتممات المنتهية، برهن
على: أن أي مجموعتين مفتوحتين غير خاليتين تتلاقيان (ومنه فإن
الفضاء ليس \omterm{def:b3:topology:hausdorff}{هاوسدورفيًا}، فليس قابلًا للتمتير)؛ وأن كل متتالية
متباينة تتقارب إلى \emph{كل} نقطة. وأين يستعمل برهان وحدانية
النهاية صفةَ هاوسدورف؟
\end{exercise}

\begin{exercise}[$\star\star$]\label{exo:b3:topology:4}
(a) برهن على أن إسقاطَي جداء $\pi_X, \pi_Y$ متصلان
و\emph{مفتوحان} (أي إن صور المفتوحات مفتوحة)، لكنهما ليسا
مغلقين في العموم (بالمثال $\{xy = 1\}$ في $\R^2$).
(b) برهن على أن متتالية في جداء قابل للعدّ $\prod_n X_n$ لفضاءات
مترية تتقارب إذا وفقط إذا تقارب كل إحداثي، وعلى أن $d(x, y) = \sum_n 2^{-n}\min(1, d_n(x_n, y_n))$ يمتّر
\omterm{def:b3:topology:constructions}{الطوبولوجيا الجدائية}.
(c) على $\R^\N$، برهن على أن «\omterm{def:b3:topology:topology}{طوبولوجيا} الصناديق» (وفيها جميع
جداءات المفتوحات مفتوحة) أدقّ تمامًا: فالمتتالية $x^{(k)} = (1/k, 1/k, \dots)$ تتقارب
إلى $0$ في \omterm{def:b3:topology:constructions}{الطوبولوجيا الجدائية} لكن ليس في \omterm{def:b3:topology:topology}{طوبولوجيا} الصناديق.
\end{exercise}

\begin{exercise}[$\star\star$]\label{exo:b3:topology:5}
الدائرة بثلاث طرائق. برهن على أن ما يلي مماثل طوبولوجيًا مثنى
مثنى، بتطبيقات صريحة: (أ) $S^1 \subseteq \R^2$؛ (ب) $\R/\Z$ (\omterm{def:b3:topology:constructions}{بطوبولوجيا
القسمة})؛ (ج) $[0,1]/(0 \sim 1)$. \emph{(من أجل (ب): برهن على أن $\R/\Z$
هاوسدورفي --- برفع صنفين إلى ممثّلين على مسافة $\leq \frac12$
--- وعلى أن $\pi([0,1])$ هو الفضاء كله، ومنه فحاصل القسمة
متراص؛ ثم استعمل \cref{cor:b3:topology:compacthomeo}.)}
\end{exercise}

\begin{exercise}[$\star\star$]\label{exo:b3:topology:6}
ليكن $X$ فضاءً متريًا.
(a) برهن على أن كل اتحاد منتهٍ لأجزاء متراصة متراص، وأن كل
تقاطع كيفي كذلك.
(b) إذا كانت $K$ متراصة و $F$ مغلقة و $K \cap F = \varnothing$، فبرهن على
$d(K, F) = \inf\{d(x,y) : x \in K, y \in F\} > 0$؛ وأعطِ مثالًا مضادًّا بمجموعتين مغلقتين منفصلتين.
(c) برهن على أن $X$ متراص إذا وفقط إذا كان \emph{كل} \omterm{def:b3:topology:continuity}{تطبيق متصل}
$f
\colon X \to \R$ محدودًا. \emph{(إذا كانت لمتتالية ما لا متتالية جزئية
متقاربة، فابنِ دالة متصلة غير محدودة محمولها قرب حدودها؛ أو
استعمل دوالًّا من نمط $x \mapsto
d(x, \cdot)$ --- وطريق نظيف واحد: إذا لم تكن
للمتتالية $(x_n)$ نقطة تراكم، فإن المجموعة $\{x_n\}$ مغلقة ومتقطّعة،
ويمتد $f(x_n) = n$ اتصالًا دون حاجة إلى تيتسه:
$f(x) = \sum_n n\,\max\bigl(0, 1 - \frac{d(x, x_n)}{r_n}\bigr)$
من أجل $r_n$ صغيرة بما يكفي.)}
\end{exercise}

\begin{exercise}[$\star\star$]\label{exo:b3:topology:7}
(عدد لوبيغ) لتكن $(U_i)$ تغطية مفتوحة لفضاء متري متراص $X$.
برهن على وجود $\delta > 0$ بحيث يقع كل جزء قطره $< \delta$ داخل
مجموعة واحدة $U_i$. \emph{(وإلا اختر $A_n$ قطرها $< 1/n$ ولا تقع
في أي $U_i$، ثم نقطة تراكم لنقط مختارة $x_n \in A_n$.)} واستنتج
مبرهنة هاين (\cref{cor:b3:topology:heineborel}(b)) من جديد.
\end{exercise}

\begin{exercise}[$\star\star$]\label{exo:b3:topology:8}
(a) برهن على أن $GL_n(\R)$ جزء مفتوح وكثيف من $M_n(\R)$، وعلى
أنه غير مترابط: فإشارة المحدد تفصله إلى قطعتين (على الأقل).
وبرهن من جهة أخرى على أن $GL_n(\C)$ مترابط بالمسارات.
\emph{(من أجل $A, B \in GL_n(\C)$: إن $\lambda \mapsto \det\bigl((1 -
\lambda)A + \lambda B\bigr)$ كثير حدود غير معدوم تطابقيًا،
ومنه فله عدد منتهٍ من الجذور في $\C$: اختر مسارًا للعناصر
$\lambda$ في $\C$ من $0$ إلى $1$ يتجنّبها.)}
(b) برهن على الكثافة: $A + \varepsilon I$ قابلة للقلب من أجل $\varepsilon > 0$ صغيرة.
\end{exercise}

\begin{exercise}[$\star\star$]\label{exo:b3:topology:9}
(a) برهن على أن المركّبات المترابطة للفضاء $\Q$ أحادية، وعلى أن
$\Q$ ليس \omterm{def:b3:topology:locallycompact}{متراصًّا محليًا} (إذ إن جوارًا \omterm{def:b3:topology:compact}{متراصًّا} للنقطة $0$ في
$\Q$ سيحتوي على $[-r, r]\cap\Q$، وغلقه في $\Q$ غير متراص:
بالقطع عند عدد أصمّ).
(b) أكمل \cref{ex:b3:topology:sinecurve}: لا يصل أي \omterm{def:b3:topology:pathconnected}{مسار} بين $(0,0)$ و $\Gamma$.
\emph{(إذا كان $\gamma = (\gamma_1, \gamma_2)$ مسارًا كهذا، فلنضع $t^* = \sup\{t : \gamma_1(t) = 0\}$؛ فمن أجل
$t
> t^*$ تتحرّك النقطة على المنحني؛ اختر $t_k \downarrow
t^*$ بحيث يصيب
$\gamma_1(t_k) \to 0$ فواصل يكون فيها $\sin$ مساويًا $\pm1$ بالتناوب --- وتوفّرها
مبرهنة القيم الوسطى --- مناقضًا اتصال $\gamma_2$ عند $t^*$.)}
\end{exercise}

\begin{exercise}[$\star\star\star$]\label{exo:b3:topology:10}
مجموعة كانتور ذات الأثلاث الوسطى $C = \bigcap_n C_n$، حيث $C_0 =
[0,1]$ وتنزع
$C_{n+1}$ الثلثَ الأوسط المفتوح من كل فترة من فترات $C_n$.
(a) برهن على $C = \bigl\{\sum_{n\geq1} a_n3^{-n} : a_n \in
\{0,2\}\bigr\}$، وعلى أن $C$ متراصة وداخلها خالٍ وليس لها
نقطة معزولة (أي إنها \emph{كاملة}).
(b) برهن على أن $C$ غير مترابطة كليًا.
(c) برهن على أن $x \mapsto (a_n(x))_n$ \omterm{def:b3:topology:continuity}{تماثل طوبولوجي} $C \to
\{0,2\}^\N$ (\omterm{def:b3:topology:constructions}{بالطوبولوجيا
الجدائية} على الفضاء المتقطّع ذي النقطتين)؛ واستنتج أن $C$ غير
قابلة للعدّ، وأن $C \times C
\cong C$.
\end{exercise}

\begin{exercise}[$\star\star\star$]\label{exo:b3:topology:11}
الإسقاط المجسّم: انطلاقًا من القطب الشمالي $N = (0, \dots,
0, 1)$ للكرة $S^n \subseteq \R^{n+1}$،
يكون التطبيق $\sigma(x) =
\frac{(x_1, \dots, x_n)}{1 - x_{n+1}}$ تماثلًا طوبولوجيًا $S^n
\setminus \{N\} \to \R^n$ (أعطِ المقلوب
صراحةً). واستنتج أن $S^n$ مماثل طوبولوجيًا للترصيص بنقطة واحدة
$\widehat{\R^n}$، وأن نزع \emph{أي} نقطة من $S^n$ يترك فضاءً
مماثلًا طوبولوجيًا للفضاء $\R^n$.
\end{exercise}

\begin{exercise}[$\star\star\star$]\label{exo:b3:topology:12}
(\omterm{ex:b3:topology:sinecurve}{منحني جيب الطوبولوجي}) لتكن
\[
T = \bigl\{(x, \sin\tfrac1x) : 0 < x \leq 1\bigr\}
\cup \bigl(\{0\}\times\intcc{-1}1\bigr) \subseteq \R^2 .
\]
(a) برهن على أن $T$ متراص، وعلى أنه \omterm{def:b3:topology:interior}{غلق} جزء المنحني.
(b) برهن على أن $T$ مترابط \emph{(فالمنحني مترابط بوصفه صورة
متصلة؛ ويبقى غلقه \omterm{def:b3:topology:connected}{مترابطًا})}.
(c) برهن على أن $T$ \emph{غير} مترابط بالمسارات: فلا يصل أي
\omterm{def:b3:topology:pathconnected}{مسار} متصل بين $(0,0)$ و $(1, \sin1)$. \emph{(إذا كان $\gamma =
(\gamma_1, \gamma_2)$
مسارًا كهذا، فلنضع $t^* = \sup\{t :
\gamma_1(t) = 0\}$؛ فبُعيد $t^*$ مباشرةً يأخذ
$\gamma_1$ جميع القيم الموجبة الصغيرة (بمبرهنة القيم الوسطى)،
ومنه يتذبذب $\gamma_2$ بين $\pm1$ على كل فترة $(t^*, t^* + \delta)$ ---
فناقض الاتصال عند $t^*$.)}
(d) اخلص إلى أن الترابط بالمسارات أقوى تمامًا من الترابط، وبرهن
على أنه لا يمكن لمثال كهذا أن يكون مفتوحًا في $\R^n$: فكل جزء
\emph{مفتوح} مترابط من $\R^n$ مترابط بالمسارات \emph{(فمجموعة
النقط التي يمكن وصلها بنقطة أساس \omterm{def:b3:topology:pathconnected}{بمسار} مفتوحة ومغلقة في
المنطلق)}.
\end{exercise}

\section{مسألة: مجموعة كانتور ومنحنٍ مالئ للمربّع}

\begin{problem}[{مسألة نهاية الأسبوع --- منحنيات بيانو موجودة،
ولماذا ليست تماثلات طوبولوجية}]\label{pb:b3:topology:1}
سنة 1890 صدم بيانو التحليلَ بواسطة \emph{تطبيق شامل متصل} $[0,1] \to [0,1]^2$:
أي منحنٍ يملأ مربّعًا. وسنبني واحدًا بأيدٍ عارية انطلاقًا من
مجموعة كانتور $C$ الواردة في \cref{exo:b3:topology:10} (ويمكن استعمال نتائجه
بحرية)، ثم نبرهن على أن مثل هذا التطبيق لا يمكن أن يكون
متباينًا: فالمربّعات ليست منحنيات. وفي كل ما يلي، تُكتب عناصر
$C$ على الصورة $x =
\sum_{n \geq 1} \frac{2b_n(x)}{3^n}$ بأرقام $b_n(x) \in
\{0, 1\}$.

\medskip
\noindent\textbf{الجزء الأول --- قراءة الأرقام قراءةً متصلة.}
\begin{enumerate}
  \item برهن على أنه إذا حقّق $x, y \in C$ الشرط $\abs{x - y} <
        3^{-m}$، فإن
        $b_n(x) = b_n(y)$ من أجل كل $n \leq m$. \emph{(فإذا كان أول رقم
        مختلف هو $b_k$، فإن $\abs{x - y} \geq \frac{2}{3^k} - \sum_{j>k}\frac2{3^j}
        = \frac1{3^k}$.)}
  \item استنتج أن كل دالة رقم $b_n \colon C \to \{0,
        1\}$ متصلة، واستنبط من جديد
        \omterm{def:b3:topology:continuity}{التماثل الطوبولوجي} $C
        \cong \{0,1\}^\N$ الوارد في \cref{exo:b3:topology:10}(c).
\end{enumerate}

\medskip
\noindent\textbf{الجزء الثاني --- تطبيق شامل متصل $C \to
[0,1]^2$.}
\begin{enumerate}[resume]
  \item نعرّف $g \colon C \to [0,1]$ بالعلاقة $g(x) = \sum_{n\geq1}
        b_n(x)\,2^{-n}$ (بقراءة أرقام كانتور على
        أنها ثنائية). برهن على أن $g$ متصل (باستعمال السؤال 1)
        وشامل. وهل هو متباين؟
  \item نعرّف $\Phi \colon C \to [0,1]^2$ بالعلاقة
        \[
        \Phi(x) = \Bigl(\ \sum_{n\geq1} b_{2n-1}(x)\,2^{-n},\
        \sum_{n\geq1} b_{2n}(x)\,2^{-n}\ \Bigr)
        \]
        (فالأرقام الفردية تعطي الفاصلة والأرقام الزوجية تعطي
        الترتيب). برهن على أن $\Phi$ متصل وشامل.
\end{enumerate}

\medskip
\noindent\textbf{الجزء الثالث --- سدّ الفجوات: منحني بيانو.}
\begin{enumerate}[resume]
  \item المتمم $[0,1] \setminus C$ اتحادٌ قابل للعدّ لفترات مفتوحة منفصلة
        $(u, v)$ (وهي الأثلاث الوسطى المنزوعة) تقع أطرافها في
        $C$. نعرّف $f \colon
        [0,1] \to [0,1]^2$ بالعلاقة $f = \Phi$ على $C$، ممدَّدًا
        \emph{تآلفيًا} على كل فجوة:
        \[
        f(t) = \frac{v - t}{v - u}\,\Phi(u) +
        \frac{t - u}{v - u}\,\Phi(v),
        \qquad t \in (u, v).
        \]
        برهن على أن $f$ معرَّف تعريفًا سليمًا وشامل على
        $[0,1]^2$.
  \item برهن على أن $f$ متصل عند كل نقطة من $[0,1]
        \setminus C$ (لأنه تآلفي
        محليًا)، وعند كل نقطة من $C$: إذا أُعطي $\varepsilon$،
        اختر $m$ يحقق $2^{-m} <
        \varepsilon/2$ واستعمل السؤال 1 للتحكم في $\Phi$
        على $C$ قرب $x$، ثم تحقق من أن الاستيفاء التآلفي لا يمكن
        أن يهرب: فعلى فجوة $(u,v) \subseteq
        (x - 3^{-m}, x + 3^{-m})$، تقع القيم $f(t)$ على القطعة
        $[\Phi(u), \Phi(v)]$، وطرفاها قريبان من $\Phi(x)$. واخلص إلى: أن $f$
        \emph{تطبيق شامل متصل} $[0,1] \to [0,1]^2$.
  \item استنتج تطبيقات شاملة متصلة $[0,1] \to [0,1]^n$ من أجل كل $n \geq 2$،
        وكذلك $\R \to \R^2$.
\end{enumerate}

\medskip
\noindent\textbf{الجزء الرابع --- لكنه ليس متباينًا أبدًا.}
\begin{enumerate}[resume]
  \item برهن على أن أي \emph{تطبيق متباين} متصل $f \colon [0,1]
        \to [0,1]^2$ سيكون
        تماثلًا طوبولوجيًا على صورته (\cref{cor:b3:topology:compacthomeo}).
  \item برهن على أن $[0,1]$ و $[0,1]^2$ غير مماثلين طوبولوجيًا:
        انزع نقطة مختارة بعناية وقارن الترابط (\cref{met:b3:topology:connectargument}).
  \item اخلص إلى: أن تطبيقًا شاملًا \omterm{def:b3:topology:continuity}{متصلًا} $[0,1] \to [0,1]^2$ لا يمكن أن
        يكون متباينًا أبدًا --- إذ إن تطبيقًا متباينًا سيجعل
        $[0,1]^2 = f([0,1])$ مماثلًا طوبولوجيًا للفترة $[0,1]$ حسب السؤال 8، وهذا
        يناقض السؤال 9. وأين دخل تراصُّ $[0,1]$ في الحجّة
        بالضبط؟
  \item (تتويج) اجمع العبرة: يوجد تطبيق شامل متصل لكن لا يوجد
        تقابل متصل $[0,1] \to
        [0,1]^2$. فماذا يقول هذا عن «البُعد» بوصفه
        مفهومًا طوبولوجيًا؟ صُغ بدقة مبرهنةً واحدة بُرهن عليها
        في هذه المسألة وعبارةً واحدة معقولة تبقى خارج أدواتنا
        (صمود المجال).
\end{enumerate}

\medskip
\noindent\textbf{الجزء الخامس --- حساب مجموعة كانتور.}
\begin{enumerate}[resume]
  \item برهن على أن $C + C = [0, 2]$: إذا أُعطي $t \in [0, 1]$،
        فاكتب $t = \sum_n d_n3^{-n}$ بأرقام $d_n \in \{0,
        1, 2\}$ واشطر كل رقم على الصورة
        $d_n = a_n + b_n$ حيث $a_n, b_n \in \{0, 1\}$؛ واخلص إلى $\frac C2 +
        \frac C2 = [0, 1]$ \emph{(وأي جزء من
        $[0,1]$ هو $\frac C2$، بدلالة الأرقام الثلاثية؟)}، ثم
        أعد التحجيم. ولا حاجة إلى أي حمل قط --- وقل لماذا يكون
        هذا هو بيت القصيد.
  \item فسّر هندسيًا: إن «غبار كانتور» $C \times
        C \subseteq \R^2$، وداخله خالٍ،
        يلقي ظلًّا \emph{كاملًا} على القطر: إذ يرسله الإسقاط
        $(x,
        y) \mapsto x + y$ على $[0, 2]$. وسجّل أيضًا التشابه الذاتي
        $C = \frac C3 \cup \bigl(\frac C3 +
        \frac23\bigr)$، وهو المعادلة الكامنة وراء كل صورة للمجموعة $C$.
  \item برهن على أن تشابك الأرقام يعرّف تماثلًا طوبولوجيًا
        $C
        \to C \times C$، واستنتج $C \cong C^n$ من أجل كل $n$ --- فمجموعة
        كانتور هي مربّعها ومكعبها و\dots{} فأي الفضاءات المألوفة
        تشاركها هذه الخاصية؟
  \item برهن على $\frac14 \in C$ \emph{(احسب نشره الثلاثي:
        $\frac14 = \sum_{k\geq1} 2\cdot3^{-2k}$)} رغم أن $\frac14$ ليس طرفًا لأي فترة منزوعة؛
        واستنتج --- بإحصاء الأطراف --- أن الأطراف تشكّل جزءًا
        قابلًا للعدّ وفعليًا من $C$.
  \item برهن على أن تطبيق القراءة الثنائية $g$ في السؤال 3 هو
        على الأكثر $2$ إلى $1$، وصِف بالضبط أي نقط $[0,1]$ لها
        سابقتان. (إذ يطوي $g$ المجموعةَ $C$ على $[0,1]$ بلصق
        عدد قابل للعدّ من الأزواج: وهو الظل التوفيقي لسلّم كانتور
        الذي نلتقيه من جديد في \cref{ch:b3:measure}.)
\end{enumerate}

\medskip
\noindent\textbf{الجزء السادس --- المجموعات المتراصة الكاملة:
كانتور في كل مكان.} نقول عن فضاء متري متراص غير خالٍ $P$ إنه
\emph{كامل} إذا لم تكن له نقطة معزولة.
\begin{enumerate}[resume]
  \item ليكن $P$ \omterm{def:b3:topology:compact}{متراصًّا} كاملًا وليكن $x \in P$ و $r > 0$. برهن
        على أن الكرة $B(x, r)$ تحتوي على نقطتين $y_0
        \neq y_1$ من $P$،
        ومنه على كرتين مغلقتين \emph{منفصلتين} $\bar B(y_0, \rho)$ و $\bar B(y_1, \rho)$
        داخل $B(x, r)$، مركزاهما نقطتان من $P$، ونصف قطرهما
        $\rho
        \leq r/4$ صغير بالقدر المطلوب. واشرح لماذا يكون الكمال
        (أي انعدام النقط المعزولة) هو بالضبط ما يسمح بتكرار هذا
        الشطر داخل \emph{كلٍّ} من الكرتين الجديدتين.
  \item كرّر: ابنِ مجموعات مغلقة $F_w \neq \varnothing$ مفهرسة بالكلمات
        الثنائية المنتهية $w$، بحيث تكون $F_{w0},
        F_{w1} \subseteq F_w$ منفصلة و
        $\operatorname{diam}F_w \leq 2^{-\abs w}
        \operatorname{diam}P$. وبرهن على أنه من أجل كل كلمة لامنتهية $\varepsilon \in \{0,1\}^\N$،
        يكون التقاطع $\bigcap_mF_{\varepsilon_1\cdots\varepsilon_m}$ نقطةً واحدة $h(\varepsilon)$
        \emph{(بخاصية التقاطع المنتهي للمتراص $P$)}.
  \item برهن على أن $h\colon\{0,1\}^\N \to P$ متباين ومتصل، واخلص إلى: \emph{أن
        كل فضاء متري متراص كامل غير قابل للعدّ} --- بل إن عدد عناصره
        لا يقل عن عدد عناصر $\{0,1\}^\N$. واستعد: أن $[0,1]$
        و $C$ غير قابلتين للعدّ.
  \item برهن على أن $h$ \omterm{def:b3:topology:continuity}{تماثل طوبولوجي} على صورته \emph{(تطبيق
        متباين متصل من متراص، \cref{cor:b3:topology:compacthomeo})}: أي إن كل فضاء متري
        متراص كامل يحتوي على نسخة مماثلة طوبولوجيًا من مجموعة
        كانتور. فمجموعة كانتور ليست شذوذًا بل هي البذرة
        الشمولية للكمال المتراص.
  \item استنتج أن لكل فضاء متري متراص قابل للعدّ نقطةً معزولة، وأعطِ
        مثالًا تكون فيه النقط المعزولة \omterm{def:b3:topology:interior}{كثيفة} لكنها ليست الفضاء
        كله: $K = \{0\} \cup
        \{\frac1n : n \geq 1\}$.
  \item (كانتور--بنديكسون من أجل $\R$) لتكن $F \subseteq \R$
        مغلقة. نقول عن $x$ إنها \emph{نقطة تكاثف} للمجموعة $F$
        إذا لاقى كل \omterm{def:b3:topology:topology}{جوار} للنقطة $x$ المجموعةَ $F$ في عدد غير
        قابل للعدّ من النقط. برهن على أن نقط التكاثف لمجموعة مغلقة
        غير قابلة للعدّ $F$ تشكّل مجموعة مغلقة كاملة غير خالية $P$،
        وأن $F
        \setminus P$ قابلة للعدّ \emph{(بتغطية النقط التي ليست نقط
        تكاثف بعدد قابل للعدّ من الفترات الناطقة التي تلاقي $F$ في
        عدد قابل للعدّ من النقط)}. واخلص إلى: أن كل جزء مغلق من $\R$
        إما قابل للعدّ وإما له عدد عناصر المتصل --- أي إن فرضية
        الاستمرار صحيحة من أجل المجموعات المغلقة.
\end{enumerate}

\medskip
\noindent\textbf{الجزء السابع --- خاتمتان: مدى الانتظام، ومدى
البعد عن الكمال.}
\begin{enumerate}[resume]
  \item (انتظام هولدر للمنحني) نضع $\alpha =
        \frac{\ln 2}{2\ln 3} \approx 0.3155$، ونلاحظ المتطابقة
        $3^\alpha = \sqrt2$. وباستعمال السؤال 1، برهن على
        \[
        \norm{\Phi(x) - \Phi(y)}_\infty \leq 2\,\abs{x -
        y}^\alpha \qquad (x, y \in C),
        \]
        ثم انشر التقدير عبر الفجوات التآلفية: برهن على أن منحني
        بيانو $f$ في السؤال 6 يحقق $\norm{f(t) - f(s)}_\infty \leq 6\abs{t - s}^\alpha$ على $[0,1]$ بأكمله
        \emph{(عالج زوجًا في الفجوة نفسها بالاستيفاء، ثم زوجًا
        عامًا بالمرور عبر النقطتين الطرفيتين من
        $C \cap [t, s]$)}.
  \item (الأُسّ $\frac12$ جدار) برهن على أنه لا يمكن لأي تطبيق
        شامل $F \colon [0,1] \to [0,1]^2$ أن يكون هولدريًا من الأُسّ $\alpha$ مع
        $\alpha > \frac12$: اقطع $[0,1]$ إلى $N$ فترةً، وحُدّ أقطار صورها،
        وأحصِ نقط الشبكة $\bigl(\frac iM, \frac jM\bigr)$، $0 \leq i, j \leq M$، التي يمكن لمجموعة
        قطرها $< \frac1M$ أن تحتويها؛ ثم اختر $M$ من الرتبة
        $N^\alpha$ واجعل $N \to \infty$. وحدّد موضع منحنينا
        ($\alpha \approx 0.3155$) بالنسبة إلى الجدار، وسجّل دون برهان أن منحني
        هيلبرت يبلغ الأُسّ الحرج $\frac12$.
  \item (المجموعات المشتقة: قياس النقص عن الكمال) من أجل $F$
        مغلقة في $\R$، لتكن $F'$ مجموعة نقط تراكم $F$ (وهي جزء
        مغلق)، ونكرّر: $F^{(0)} = F$، $F^{(i+1)}
        = (F^{(i)})'$. تحقق من أن
        \[
        K = \{0\} \cup \Bigl\{\frac1m\Bigr\}_{m \geq 1} \cup
        \Bigl\{\frac1m + \frac1n : m \geq 1,\ n > m^2\Bigr\}
        \]
        مجموعة متراصة قابلة للعدّ تحقق $K' = \{0\} \cup
        \{\frac1m\}$ و $K'' = \{0\}$
        و $K''' = \varnothing$ \emph{(تحقق من أن العنقود النوني من الرتبة $m$
        يعيش في الفترة $(\frac1m, \frac1{m-1})$، ومنه فالعناقيد لا تتشابك)}، ومن
        أن نقطها المعزولة \omterm{def:b3:topology:interior}{كثيفة} في $K$، كما يتنبّأ السؤال 21.
        وصِف التراجع الذي يُنتج، من أجل كل $j$، مجموعة متراصة
        قابلة للعدّ $K_j$ تحقق $K_j^{(j)} = \{0\}$ و $K_j^{(j+1)} =
        \varnothing$، وقارن ذلك بالمجموعات
        الكاملة في السؤال 22، التي لا يحرّكها الاشتقاق أبدًا:
        فالرتبة المنتهية هي النقيض التام للكمال.
\end{enumerate}
\end{problem}
